\documentclass[11pt]{article}
\usepackage{amsmath,amssymb,mathtools,bm}
\usepackage[margin=1in]{geometry}

\title{A Higher-Order LS2010 Flux Kernel for the N-D Acausal FIF Case}
\author{GPT-5.4 Pro Extended, Thomas D. DeWitt}
\date{}

\begin{document}
\maketitle

\section{Goal}

The aim is to build the N-D analogue of the higher-order 1D acausal flux kernel:
a drop-in replacement for the \texttt{FIF\_ND} flux-kernel call that stays in the
Lovejoy--Schertzer (2010) real-space framework, keeps the same large-scale power-law
tail, but removes several discrete small-scale bias modes instead of only the single
mode corrected in the present implementation.

As in the 1D file, this is a five-mode LS2010-style correction rather than a literal
closed-form fifth-order symbolic Euler--Maclaurin cancellation.

\section{What the current N-D code is doing}

Your current N-D FIF implementation uses the flux-kernel parameters
\begin{equation}
\alpha' = \frac{1}{1-1/\alpha},
\qquad
\texttt{flux\_exponent} = -\frac{d_{\mathrm{eff}}}{\alpha'},
\qquad
\texttt{flux\_norm\_ratio\_exp} = -\frac{d_{\mathrm{eff}}}{\alpha},
\qquad
\texttt{flux\_final\_power} = \frac{1}{\alpha-1},
\end{equation}
where $d_{\mathrm{eff}}$ is either the spatial dimension or the user-supplied
\texttt{scale\_metric\_dim}.  After the final power is applied, the flux kernel has the
usual tail
\begin{equation}
 g_0(\bm x) \propto \|\bm x\|^{-d_{\mathrm{eff}}/\alpha}.
\end{equation}
This is exactly how the current \texttt{FIF\_ND} and \texttt{create\_kernel\_LS2010}
code are wired. 

\section{Why the same LS2010 logic extends to higher dimensions}

LS2010 explicitly state that the discrete normalization generalizes from
\begin{equation}
N_{D,f} = \frac{N_D}{\Delta X^D}
\end{equation}
to arbitrary $D$, and in Appendix~B they show that for isotropic higher-dimensional
systems the angular integration leaves the same radial singular structure as in 1D.
In their notation, after integrating over the angles on the displaced
$r=\mathrm{const}$ hypersurfaces, one recovers the same logarithmic core measure
\begin{equation}
\frac{dr}{r(1-r^2)}\,\Omega_D,
\end{equation}
with $N_D=\Omega_D$.  So the higher-dimensional isotropic problem reduces to the
same radial structure as the 1D calculation; this is why they describe the extension
as straightforward.

This strongly suggests that the finite-size correction should remain a \emph{radial}
small-$r$ modification of the kernel, now with the $d_{\mathrm{eff}}$-dimensional
power-law tail.

\section{Leading discrete bias in N dimensions}

In 1D, LS2010 identify the dominant bad term as the discrete
\(S_f^{(1,1)}\) contribution.  The same mechanism persists in higher dimensions:
it comes from the neighborhoods of the two singularities at
\(\bm x=\bm 0\) and \(\bm x=\bm \Delta\).

For an isotropic N-D flux kernel with tail
\begin{equation}
 g_0(r) \sim r^{-d_{\mathrm{eff}}/\alpha},
\end{equation}
we can estimate the higher-dimensional analogue of the LS2010 term by combining
one far-field factor $g_0(\Delta)$ with a local integral of $g_0^{\alpha-1}$:
\begin{equation}
S_f^{(1,1)}(\Delta)
\sim
2\alpha\, g_0(\Delta)
\int_1^{\Delta/2}
\rho^{d_{\mathrm{eff}}-1}
\bigl(g_0(\rho)\bigr)^{\alpha-1}
\,d\rho.
\end{equation}
Since
\begin{equation}
\bigl(g_0(\rho)\bigr)^{\alpha-1}
\sim
\rho^{-d_{\mathrm{eff}}(\alpha-1)/\alpha},
\end{equation}
the integrand is
\begin{equation}
\rho^{d_{\mathrm{eff}}-1-d_{\mathrm{eff}}(\alpha-1)/\alpha}
=
\rho^{d_{\mathrm{eff}}/\alpha - 1},
\end{equation}
so the radial integral contributes
\begin{equation}
\int_1^{\Delta/2} \rho^{d_{\mathrm{eff}}/\alpha -1} d\rho
=
A + B\,\Delta^{d_{\mathrm{eff}}/\alpha}.
\end{equation}
Multiplying by
\(g_0(\Delta)\sim \Delta^{-d_{\mathrm{eff}}/\alpha}\)
therefore yields
\begin{equation}
S_f^{(1,1)}(\Delta)
=
C_0 + C_1\,\Delta^{-d_{\mathrm{eff}}/\alpha} + \cdots .
\end{equation}
So the higher-dimensional leading non-logarithmic mode is the natural N-D analogue
of the 1D LS2010 correction, namely
\begin{equation}
\Delta^{-d_{\mathrm{eff}}/\alpha}.
\end{equation}
This also explains why the small-scale bias is already milder in 2D than in 1D:
for fixed $\alpha$, the correction falls faster when $d_{\mathrm{eff}}$ increases.

\section{Five-mode N-D LS2010 ansatz}

The direct N-D extension of the 1D higher-order kernel is the radial ansatz
\begin{equation}
 g_M(r)
 =
 r^{-d_{\mathrm{eff}}/\alpha}
 \left(
   1 + \sum_{m=1}^{M} c_m e^{-r/\lambda_m}
 \right)^{1/(\alpha-1)},
 \qquad M=5.
 \label{eq:nd_ansatz}
\end{equation}
This keeps exactly the same LS2010 philosophy:
\begin{itemize}
\item the kernel is modified only at small $r$,
\item the far-field tail is unchanged,
\item the correction remains strictly positive provided
\(1+\sum_m c_m e^{-r/\lambda_m}>0\).
\end{itemize}

The key simplification is the same as in 1D:
\begin{equation}
 g_M(r)^{\alpha-1}
 =
 r^{-d_{\mathrm{eff}}(\alpha-1)/\alpha}
 \left(
   1 + \sum_{m=1}^{M} c_m e^{-r/\lambda_m}
 \right),
\end{equation}
so the correction enters linearly in the quantity that drives the LS2010 bias.
At the same time,
\begin{equation}
 g_M(r)^\alpha
 =
 r^{-d_{\mathrm{eff}}}
 \left[1 + O\!igl(e^{-r/\lambda_1}\bigr)\right],
\end{equation}
so the one-point logarithmic scaling is unchanged; only the additive finite-size
constants are altered, which is entirely in the LS2010 spirit.

\section{Discrete fitting criterion in N dimensions}

For the N-D odd grid, let
\begin{equation}
\mathcal O_{\bm L}
=
\prod_{j=1}^{D}
\{-L_j+1,-L_j+3,\dots,L_j-1\}
\end{equation}
be the LS2010 odd lattice.  For a displacement vector \(\bm \Delta\) with even
components, define the exact discrete LS2010 quantity
\begin{equation}
S_f(\bm \Delta; \bm c)
=
\sum_{\bm i \in \mathcal O_{\bm L}}
\left[g_M(\bm i-\bm \Delta) + g_M(\bm i)\right]^\alpha.
\end{equation}
For a perfectly scaling isotropic kernel, this should be affine in
\(\log \|\bm \Delta\|\).

In practice the implementation evaluates \(S_f\) along a small set of rays
\begin{equation}
\bm \Delta = 2k\,\hat{\bm d}_r,
\qquad
k=1,\dots,K,
\end{equation}
where the default choices are the coordinate axes and, for the default Euclidean
metric, also the full diagonal.  The factor of 2 is the physical LS2010 odd-grid
spacing; in the code, the shifts are represented in index units, so the calibration
uses the equivalent lags
\(k\,\|\hat{\bm d}_r\|\).

Collect all sampled values in a vector
\begin{equation}
\bm S(\bm c)
=
\bigl(S_f(\bm \Delta_1;\bm c),\dots,S_f(\bm \Delta_N;\bm c)\bigr)^T,
\end{equation}
and form the projector orthogonal to pure logarithmic scaling,
\begin{equation}
P_\perp = I - X(X^T X)^{-1}X^T,
\qquad
X=
\begin{bmatrix}
1 & \log \ell_1\\
\vdots & \vdots\\
1 & \log \ell_N
\end{bmatrix},
\end{equation}
where \(\ell_n\) is the sampled lag magnitude.  The fitted coefficients are then
obtained by minimizing
\begin{equation}
\Phi(\bm c)
=
\|W P_\perp \bm S(\bm c)\|_2^2 + \lambda\|\bm c\|_2^2,
\end{equation}
with optional small-lag weights \(W\) and a tiny Tikhonov regularizer \(\lambda\).

This is the N-D analogue of the 1D construction: instead of cancelling just the
first asymptotic term, the kernel is tuned so the full discrete N-D
\(S_f\)-surface is as log-affine as possible over the range where the bias is most
visible.

\section{Analytic Jacobian}

The Gauss--Newton step uses the exact derivatives
\begin{equation}
\frac{\partial g_M(r)}{\partial c_m}
=
\frac{1}{\alpha-1}
\,g_M(r)
\frac{e^{-r/\lambda_m}}
     {1+\sum_n c_n e^{-r/\lambda_n}}.
\end{equation}
Hence
\begin{equation}
\frac{\partial S_f(\bm \Delta)}{\partial c_m}
=
\alpha
\sum_{\bm i \in \mathcal O_{\bm L}}
\left[g_M(\bm i-\bm \Delta)+g_M(\bm i)\right]^{\alpha-1}
\left(
\frac{\partial g_M(\bm i-\bm \Delta)}{\partial c_m}
+
\frac{\partial g_M(\bm i)}{\partial c_m}
\right).
\end{equation}
After projection and weighting, the normal equations are
\begin{equation}
\left(J^T J + \lambda I\right)\delta \bm c = -J^T \bm r,
\qquad
\bm r = W P_\perp \bm S(\bm c),
\end{equation}
with a short damped line search to keep
\(1+\sum_m c_m e^{-r/\lambda_m}>0\)
everywhere.

\section{Handling custom scale metrics}

Your N-D FIF code already supports a user-supplied \texttt{scale\_metric} and a
possibly non-integer \texttt{scale\_metric\_dim}.  The new kernel preserves that
interface.

The only change is that the Euclidean odd-grid radius $r$ in
Eq.~\eqref{eq:nd_ansatz} is replaced by the supplied metric field, and the
Euclidean dimension $D$ is replaced by
\(d_{\mathrm{eff}} = \texttt{scale\_metric\_dim}\).
So the implementation is still radial with respect to the chosen metric.

For the default isotropic metric, the fit uses axes plus the full diagonal.  For a
custom metric, the default fit uses axes only; this keeps the internal calibration
agnostic about the exact angular structure of the metric while still removing the
worst small-scale bias modes.

\section{Code mapping}

The code file defines two public entry points.

\begin{enumerate}
\item
\texttt{create\_kernel\_flux\_LS2010\_highorder\_nd(size, alpha, ...)}\\
This is the direct N-D kernel builder.

\item
\texttt{create\_kernel\_LS2010\_flux\_highorder\_nd\_from\_params(...)}\\
This is the thin wrapper intended for the existing \texttt{FIF\_ND} call site.
The swap is simply
\begin{verbatim}
kernel1 = create_kernel_LS2010_flux_highorder_nd_from_params(
    sim_size, flux_exponent, flux_norm_ratio_exp,
    causal=False, outer_scale=outer_scale,
    outer_scale_width_factor=outer_scale_width_factor,
    final_power=flux_final_power,
    scale_metric=scale_metric,
    scale_metric_dim=scale_metric_dim,
    order=5,
)
\end{verbatim}
in place of the current
\begin{verbatim}
kernel1 = create_kernel_LS2010(
    sim_size, flux_exponent, flux_norm_ratio_exp,
    causal=False, outer_scale=outer_scale,
    outer_scale_width_factor=outer_scale_width_factor,
    final_power=flux_final_power,
    scale_metric=scale_metric,
)
\end{verbatim}
\end{enumerate}

\section{What to expect}

This kernel should help most in the regime where the current N-D LS2010-style
single correction is still visibly too smooth at the smallest resolved scales.  The
benefit should be strongest in 2D/3D acausal flux generation with \(\alpha>1\),
while still preserving the same real-space extremal-L\'evy convolution workflow,
odd-grid LS2010 construction, and outer-scale tapering used in your current code.

\end{document}
